\section{Current Experiments}
\label{sec:experiments}

\subsection{Reinforcement Learning}
\begin{itemize}
  \item In reinforcement learning (RL), an agent learns by trial and error,
        receiving rewards for good outcomes~\cite{sutton2018}.
  \item Landmark results: human-level play on Atari games~\cite{mnih2015},
        \emph{Dota~2} with OpenAI Five~\cite{berner2019}, and Grandmaster
        level in \emph{StarCraft~II} with AlphaStar~\cite{vinyals2019}.
  \item Important distinction: these systems aim to \emph{win}, while NPCs
        should \emph{entertain}; strong play is not the same as good game
        design.
  \item Toolkits such as Unity ML-Agents make RL accessible to game
        developers~\cite{juliani2018}.
  \item Barriers to commercial use: long training times, unpredictable
        behavior, and the need to retrain when the game changes.
  \item Practical uses today: automated playtesting, balancing, and
        training opponents for racing or fighting games.
\end{itemize}

\subsection{Neuroevolution and Learning from Players}
\begin{itemize}
  \item Neuroevolution evolves neural networks over many generations.
        \emph{NERO} made training an army of NPCs the core
        gameplay~\cite{stanley2005}.
  \item Imitation learning copies human play to create more human-like
        bots~\cite{yannakakis2018}.
  \item Player modeling lets NPCs and games adapt to individual play
        styles~\cite{yannakakis2018}.
\end{itemize}

\subsection{Large Language Models as NPCs}
\begin{itemize}
  \item LLMs allow NPCs to hold open-ended conversations instead of picking
        from fixed dialogue options~\cite{gallotta2024}.
  \item \emph{Generative Agents} placed 25 LLM-driven characters in a small
        town; they planned their days, formed relationships and even
        organized a party on their own~\cite{park2023}.
  \item Architecture explained simply: memory of past events, reflection
        to form higher-level thoughts, and planning of daily
        activities~\cite{park2023}.
  \item \emph{Voyager} used an LLM to explore and learn new skills in
        \emph{Minecraft} without human help~\cite{wang2023}.
  \item Gallotta et al.\ survey the many roles LLMs can play in games,
        including as NPCs, designers and players~\cite{gallotta2024}.
  \item The idea links back to earlier interactive drama
        research~\cite{mateas2005,evans2014}, now with far more flexible
        language.
\end{itemize}

\subsection{Limitations of LLM-Driven NPCs}
\begin{itemize}
  \item \emph{Hallucination}: the model may invent facts that contradict
        the game world or its story~\cite{ji2023}.
  \item Consistency: characters can drift out of role or be manipulated by
        players into saying things the designer never intended.
  \item Cost and latency: each response requires heavy computation, which
        clashes with tight frame budgets (Section~\ref{sec:tradeoffs}).
  \item Designer control is reduced, the same trade-off seen with planning
        and RL, but stronger.
  \item Ethical concerns: bias in training data, harmful content, and the
        environmental cost of large models~\cite{bender2021}.
\end{itemize}

\subsection{Emerging Hybrid Approaches}
\begin{itemize}
  \item Promising direction: keep classical systems (behavior trees,
        planners) in charge of actions, and use LLMs only for dialogue or
        high-level intent~\cite{gallotta2024}.
  \item Guardrails restrict what the model may say or do in the game
        world.
  \item Smaller local models could reduce cost and latency.
\end{itemize}

\subsection{Evaluating Experimental NPCs}
\begin{itemize}
  \item Win rates do not capture what matters for NPCs; player studies and
        believability ratings are needed~\cite{warpefelt2017,hingston2009}.
  \item Player enjoyment frameworks such as GameFlow offer structured
        criteria~\cite{sweetser2005}.
  \item Few studies so far compare LLM-based and classical NPCs with real
        players, which is an open research gap.
\end{itemize}
